\documentclass[10pt]{article}

\title{Damage Severity Estimation from a Single Product Image}
\author{cv-research-lab}
\date{2026-09-22}

\begin{document}
\maketitle

\begin{abstract}
Shell created on day 1. The abstract will later state the three-level definition, the comparison result, and where the model fails. This file does not contain measured numbers.
\end{abstract}

\noindent\textbf{Keywords:} damage severity, ordinal classification, visual inspection

\section{Introduction}
The task is to assign one ordered grade, none / minor / severe, when the product category is already known.
The three claims are written in \texttt{paper/contributions.md}.
They are: a written three-level definition, a comparison that includes locate-then-grade, and a statement of when the model fails.

\section{Related Work}
Subsections still empty: surface defects, ordinal classification, condition assessment.

\section{Method}
\subsection{Label definition}
\subsection{Whole-image classifier and ordinal loss}
\subsection{Locate then grade}

\section{Experiments}
\subsection{Data and split}
\subsection{Metrics and implementation}
\subsection{Main comparison}
\subsection{Locate-then-grade}
\subsection{Ablation and ordinal error}
\subsection{Failure cases and real photos}

\section{Conclusion}
Limitations wait for the frozen tables: reflection, cross-product shift, annotation subjectivity, and the gap between public labels and real condition grades.

\end{document}
